%!TEX root = ../csuthesis_main.tex
% 设置中文摘要
\keywordscn{持续学习；医学图像分类；语义分割；时间序列分类；多模态大语言模型}
%\categorycn{TP391}
\begin{abstractzh}

持续学习需要解决问题是：当数据分批流入并且类别的集合不断变化的情况下，模型既要适应新的任务，又不能忘记之前学到的知识。当前大多数的工作都是利用梯度下降进行在线更新，容易受到梯度冲突的影响导致发生灾难性遗忘；而重放样例的方法虽然可行但是带来了很大的空间消耗以及隐私问题。本文另辟蹊径，从闭式解的角度出发，对基于解析学习的持续学习方法进行了总结归纳。

研究从整理递归岭回归闭式解开始，将其归纳为一个整体的数学表达式：在起始部分训练出编码器固定不动，只对解析分类头进行递归更新，使得增量学习可以转化为有正则化的最小二乘法的问题；同时给出了一种不需要使用之前所有原始样本的递推公式并且进行了证明，在此基础上本文在几个不同任务中进行算法实现及实验测试。其中包括：医学图像疾病分类用于说明解析持续学习学习在隐私安全的应用场景；时间序列分类结合多尺度特征融合以及随机高维映射解决灾难性遗忘和类内差异的问题；对于类别增量语义分割提出一种将随机高维映射和伪标签结合起来的持续学习分割方案，可以同时处理二维图像和三维点云；对于多模态语言模型提出一种基于解析式低秩专家选择头来完成高效的专家路由的方法。

综合全面的实验显示所提出的方法对于语义分割、医学图像分类、时间序列分类以及多模态语言模型专家路由这几种任务均能获得优秀的的结果,同时大大减少持续学习期间的学习成本并且无需使用任何以往的数据进行回放。

\end{abstractzh}